\section{Competencias específicas}
En conformidad con la \textbf{\emph{Memoria de Verificación del Título de Grado en Ingeniería Informática}} \cite{ulpgc_memoria2019}, se exponen a continuación las competencias específicas relacionadas con este proyecto.

\subsection{Competencias comunes a la rama informática}
\begin{itemize}
    \item \textbf{CI1~-~\emph{Capacidad para diseñar, desarrollar, seleccionar y evaluar aplicaciones y sistemas informáticos, asegurando su fiabilidad, seguridad y calidad, conforme a principios éticos y a la legislación y normativa vigente.}}

    Esta competencia se aborda en varios capítulos del proyecto, como;
    estado del arte, normativa y legislación, y tecnologías y herramientas.

    \item \textbf{CI2~-~\emph{Capacidad para planificar, concebir, desplegar y dirigir proyectos, servicios y sistemas informáticos en todos los ámbitos, liderando su puesta en marcha y su mejora continua y valorando su impacto económico y social.}}

    Se trabaja esta competencia durante todo el proyecto, comenzando por la selección de la metodología a utilizar, pasando por la planificación y terminando con la ejecución del proyecto.

    \item \textbf{CI3~-~\emph{Capacidad para comprender la importancia de la negociación, los hábitos de trabajo efectivos, el liderazgo y las habilidades de comunicación en todos los entornos de desarrollo de software.}}

    Podemos ver esta competencia reflejada tanto en el capítulo de estado del arte como en el desarrollo. Por un lado, a través de la investigación y análisis del contexto y alcance del proyecto, y por otro lado, a través del análisis de requisitos y diseño.
    
    \item \textbf{CI6~-~\emph{Conocimiento y aplicación de los procedimientos algorítmicos básicos de las tecnologías informáticas para diseñar soluciones a problemas, analizando la idoneidad y complejidad de los algoritmos propuestos.}}

    Se cubre esta competencia en el capítulo de desarrollo, específicamente en la construcción de un algoritmo que permite el cálculo de acordes.

    \item \textbf{CI7~-~\emph{Conocimiento, diseño y utilización de forma eficiente de los tipos y estructuras de datos más adecuados a la resolución de un problema.}}

    Queda demostrada en el capítulo de desarrollo, donde se diseñan e implementan estructuras de datos adaptadas a las características particulares de la funcionalidad en cuestión.

    \item \textbf{CI8~-~\emph{Capacidad para analizar, diseñar, construir y mantener aplicaciones de forma robusta, segura y eficiente, eligiendo el paradigma y los lenguajes de programación más adecuados.}}

    Esta competencia se desarrolla en el estado del arte, tecnologías y herramientas y desarrollo, donde se ejecutan los aspectos de análisis, selección de herramientas, diseño e implementación.

    \item \textbf{CI12~-~\emph{Conocimiento y aplicación de las características, funcionalidades y estructura de las bases de datos, que permitan su adecuado uso, y el diseño y el análisis e implementación de aplicaciones basadas en ellos.}}

    En el capítulo de desarrollo se expone el conocimiento, diseño e implementación de funcionalidades y bases de datos, vitales para la entrega de valor de la aplicación web al usuario.
    
    \item \textbf{CI16~-~\emph{Conocimiento y aplicación de los principios, metodologías y ciclos de vida de la ingeniería del software.}}

    Se justifica esta competencia a través de los capítulos de metodología y planificación, y desarrollo. A través de los anteriores, se selecciona y explica la metodología y planificación sobre la que se desarrollará el proyecto, además de la demostración de su correcta ejecución.
    
    \item \textbf{CI17~-~\emph{Capacidad para diseñar y evaluar interfaces persona computador que garanticen la accesibilidad y usabilidad a los sistemas, servicios y aplicaciones informáticas.}}

    Esta competencia se expone en los capítulos de estado del arte y desarrollo, evaluando las soluciones ya existentes y planteando nuevas propuestas para el diseño de experiencias e interfaces de usuario del proyecto.

    \item \textbf{CI18~-~\emph{Conocimiento de la normativa y la regulación de la informática en los ámbitos nacional, europeo e internacional.}}

    Se dedica un capítulo exclusivamente a la exposición de la normativa y legislación, además de una sección sobre licencias en el capítulo de tecnologías y herramientas.
    
\end{itemize}

\subsection{Competencias de tecnología específica}
\begin{itemize}    
    \item \textbf{TI6~-~\emph{Capacidad de concebir sistemas, aplicaciones y servicios basados en tecnologías de red, incluyendo Internet, web, comercio electrónico, multimedia, servicios interactivos y computación móvil.}}

    Esta competencia se materializa en el resultado final del proyecto, con la existencia de una aplicación web funcional.
    
\end{itemize}

\subsection{Competencia de trabajo de fin de grado}
\begin{itemize}
    \item \textbf{TFG~-~\emph{Ejercicio original a realizar individualmente y presentar y defender ante un tribunal universitario, consistente en un proyecto en el ámbito de las tecnologías específicas de la Ingeniería en Informática de naturaleza profesional en el que se sinteticen e integren las competencias adquiridas en las enseñanzas.}}

    Esta competencia se cumple con la realización integral del proyecto, que incluye tanto la aplicación desarrollada como el presente documento, reflejando la capacidad de síntesis e integración de las competencias adquiridas a lo largo del grado.

\end{itemize}
